\section{Mission, périmètre et limites (G1)}

La direction générale nous a mandatés pour confronter les déclarations, les documents organisationnels et les extraits techniques de MEDICA-SERVICES, afin d'établir le niveau de rigueur souhaitable et la maîtrise réellement démontrée des processus SSI. L'audit porte exclusivement sur les pièces \piece{P01} à \piece{P17} et \piece{T01} à \piece{T10}. Chaque conclusion distingue le fait documenté, la déclaration, l'interprétation de l'équipe et l'information manquante.

La méthode recommande un périmètre homogène en termes de responsabilité. Nous retenons le service d'hébergement et d'exploitation que MEDICA-SERVICES rend à ses clients, placé sous l'autorité de la direction générale. Ce choix couvre les actifs qui portent les données de patients et laisse hors champ les systèmes propres aux établissements.

\begin{table}[h]
\centering\small
\caption{Périmètre retenu}
\begin{tabularx}{\linewidth}{@{}L{3.2cm} Y L{2.4cm}@{}}
\entete \hd{Élément} & \hd{Contenu retenu} & \hd{Pièces} \\
Autorité et responsabilités & DG, autorité responsable qui approuve la stratégie. DSI de 14 personnes chargée de l'application. RSSI à mi-temps rattaché à la DSI. DPO externalisé & P02, P09 \\
Bénéficiaires & 12 établissements de soins privés, reprise contractuelle sous 8 h des applications critiques & P01 \\
Actifs essentiels & MEDI-PORTAL et sa base de données, identifiants patients et informations de santé, ERP, ticketing & P03, P05, T01 \\
Services supports & 35 serveurs ou VM sous Windows et Linux, Active Directory, Microsoft 365, VPN et bastion d'administration, sauvegardes, deux salles dont une de secours chez un prestataire & P03, T01, T09 \\
Interconnexions & Portail HTTPS publié sur Internet, liaisons contractuelles avec les clients, réplication vers le site de secours, SaaS tiers & P03, P07, T01 \\
Limites et exclusions & SI des établissements clients, fonctionnement interne des SaaS tiers et du prestataire de secours, réseaux des agences de Bordeaux et Montpellier non décrits, aucun test actif & P01, T01, T04 \\
\end{tabularx}
\end{table}

L'audit se fonde sur des fragments. Les extraits de configuration ne sont pas accompagnés des journaux de trafic, l'inventaire des dépendances applicatives est partiel et les retards de correctifs sont déclaratifs \piece{(T01, T02, T06)}. Une absence de justificatif est donc traitée comme une incertitude à lever et non comme la preuve qu'une action n'a jamais eu lieu. Trois informations manquantes pèsent particulièrement sur nos conclusions et justifient des demandes de preuves précises.

\begin{table}[h]
\centering\small
\caption{Informations manquantes et demandes de preuves}
\begin{tabularx}{\linewidth}{@{}L{4.1cm} Y Y@{}}
\entete \hd{Information manquante} & \hd{Conséquence pour l'audit} & \hd{Demande de preuve ou question d'entretien} \\
Objectifs de reprise par service (RTO, RPO) et dépendances applicatives \piece{(T01, T09)} & Impossible de vérifier que l'engagement de 8 h est tenable ni de hiérarchiser les tests de restauration & Analyse d'impact par service, annexes contractuelles de reprise, convention avec le site de secours. Entretien avec le responsable infrastructures \\
Liste exhaustive des comptes privilégiés, y compris SaaS, et politique d'accès conditionnel \piece{(T04)} & La couverture MFA de 60 \% ne vaut que pour l'export fourni & Export complet de l'annuaire et des consoles SaaS, règles d'accès conditionnel, registre des exceptions approuvées \\
Suite donnée à l'alerte AL-1048 et intégrité des journaux \piece{(T07, T08)} & Impossible de qualifier l'événement du 9 octobre & Journaux bruts horodatés et scellés, historique du ticketing, entretien contradictoire avec le titulaire d'admin-old et l'astreinte de la nuit \\
\end{tabularx}
\end{table}
